\paragraph{Réponse : en partie seulement.} Évaluées avec le métamodèle v0.8, les compositions optimales de la v0.7
pour S1, S2 et S3 violent leur contrainte de dette : $+9$ points en 2047 pour S1, $+6$ points en 2037 pour S2 et $+23$
points en 2042 pour S3. Elles reposaient sur les gains budgétaires que la fermeture de la v0.7 attribuait aux mesures
d'offre et d'activité. Les compositions de S4 et S5 restent admissibles avec la distribution statique ; avec la
distribution après propagation, elles violent la contrainte de perte maximale par décile (le premier décile perd
4,1 \% dans S4 et 2,3 \% dans S5), parce que la baisse de 25 \% des allocations chômage qu'elles contiennent fait
désormais apparaître ses perdants.

\paragraph{Ce qui change dans les compositions réoptimisées.} Avec la distribution statique, la v0.8 garde la
structure des solutions v0.7 (quatre à sept leviers communs sur huit) mais réduit leur ambition budgétaire. Avec la
distribution après propagation, deux changements nets apparaissent : la baisse des allocations chômage disparaît des
stratégies S3, S4 et S5 (elle n'est plus retenue que dans 1 \% des tirages de préférences, contre 33 \% en v0.7), et S1
abaisse l'âge légal à 62 ans (retenu dans 54 \% des tirages), ce qui augmente le niveau de vie des trois premiers
déciles mais coûte 2,7 \% de PIB en 2047. Le front de Pareto est plus large : 70 \% des portefeuilles aléatoires
étaient dominés en v0.7, 30 \% seulement avec la distribution dynamique, parce que les critères distributifs entrent
davantage en conflit avec les critères d'activité et de finances publiques.
